% Figure: idealised Langmuir-probe I-V characteristic, showing the three
% regions: ion saturation, the exponential electron-retardation region
% whose slope gives T_e, and electron saturation. Data computed inline in
% code/langmuir_probe.py and saved to data/langmuir_iv.csv.
\begin{figure}[htbp]
\centering
\begin{tikzpicture}
\begin{axis}[
  width=0.78\textwidth, height=0.5\textwidth,
  xlabel={probe bias $V$ (V)},
  ylabel={probe current $I$ (mA)},
  xmin=-40, xmax=25, ymin=-3, ymax=12,
  grid=both, grid style={enggray!20},
  every axis plot/.append style={thick},
]
  \addplot[engblue, mark=*, mark size=0.7pt] table[
    col sep=comma, x=bias_V, y=current_mA]
    {volumes/04-energy/ch13-plasma-physics/data/langmuir_iv.csv};
  \draw[dashed, enggray] (axis cs:-40,0) -- (axis cs:25,0);
  \node[engred, font=\scriptsize, anchor=west] at (axis cs:-38,-1.6)
    {ion saturation};
  \node[font=\scriptsize, anchor=south] at (axis cs:-8,2.5)
    {exponential: slope $\to T_e$};
  \node[engred, font=\scriptsize, anchor=east] at (axis cs:24,10.4)
    {electron saturation};
  \draw[dashed, enggray] (axis cs:5,-3) -- (axis cs:5,12);
  \node[enggray, font=\scriptsize, anchor=south west] at (axis cs:5.3,8.5)
    {$V_p$};
\end{axis}
\end{tikzpicture}
\caption[Langmuir-probe characteristic]{Idealised current-voltage
characteristic of a Langmuir probe. Strongly negative bias repels all
electrons and the probe draws only the small ion-saturation current.
Raising the bias lets a growing fraction of the electron distribution
reach the probe; in this electron-retardation region the current rises
exponentially, and the slope of $\ln(I-I_{\mathrm{ion}})$ against $V$
gives the electron temperature $T_e$. Above the plasma potential $V_p$
the probe collects the full electron flux and the current saturates;
the saturation values fix the density $n_e$. The data are generated by
\texttt{code/langmuir\_probe.py}.}
\label{fig:vol04ch13:langmuir-iv}
\end{figure}
